\subsection{The actual terminal degree extraction}
\label{actual-terminal-degree}

Retain the original place, normalized tower and compact coefficient
class, and write $\mathcal V_H(a)$ for their orbital vector.  The
Hermitian form is linear in its first variable:
$h(v,w)=w^*Jv$.  Put $R=r+1$, $d=2r+k+1$, and let $D_j$ be the
prescribed leading paired determinant of size $j$.  For real $y$ set
\[
 \mathcal A_yH=D_R^{-1}\operatorname{AA}_r
                 \left(y+\frac{2r+k}{2}\right)H.
\]
Here $\operatorname{AA}_r$ is the original contracted column minor,
with its coefficients evaluated at the point of differentiation;
its ordered factors have parameters $u,u-1,\ldots,u-r+1$.
The definition is extended by zero wherever the original inverse is
zero.  A compact coefficient with positive Gram cutoff vanishes in
a neighborhood of that set, so this convention causes no loss of
smoothness.

\begin{theorem}\label{terminal-degree-actual}
Let $a=h\mathbin{\|}(s,t)$ satisfy the original terminal head conditions,
with $s\ge2$, $p=|h|$, $p+s=R$, and $t\ne0$.  For every original
compact coefficient $H$, the orbital vector
$y\mapsto\mathcal V_{\mathcal A_yH}(a)$ is a polynomial of degree
at most $p$, whose coefficient of $y^p$ is
\begin{equation}\label{terminal-degree-leading}
             \lambda\mathcal V_H(a),
             \qquad \lambda=\frac{(-1)^{s-1}}{t}.
\end{equation}

Keep also the prescribed operator $\Theta$, its actual regular
radial identity, and its original differential realization
\texttt{InDDplus2}.  Let $w$ be any old word and $e$ any of its
tensor indices.  For any $K\ge0$, take arbitrary injective real
grids $\{y_{a,0},\ldots,y_{a,p}\}$, $0\le a<K$, such that every
word obtained by prepending the $K$ chosen plus letters is admissible
for the original cut.  Put
\[
       \ell_{a,n}=\prod_{\substack{0\le m\le p\\m\ne n}}
                             (y_{a,n}-y_{a,m})^{-1}.
\]
Then, with the prescribed tail-first word order,
\begin{align}\label{terminal-degree-grid}
 \lambda^K\mathcal V_{\Theta_{w,e}F}(a)
 =\sum_{n\in\{0,\ldots,p\}^K}
       \left(\prod_{b=0}^{K-1}\ell_{b,n_b}\right)
       \mathcal V_{\Theta_{(+y_{0,n_0},\ldots,+y_{K-1,n_{K-1}},w),
                                      (0,\ldots,0,e)}F}(a).
\end{align}
Thus the original \texttt{PolyTop}, \texttt{TopFirst} and
\texttt{DegTop} obligations hold, with all their head, old-word,
component and grid quantifiers.  This does not construct $\Theta$;
its two stated source properties remain hypotheses of this theorem.
\end{theorem}

\begin{proof}
The original orbital vector is zero for an invalid head, for a
head of nonpositive Gram, or when no actual unitary representative
exists.  These conventions are independent of the coefficient.
Thus the polynomial assertion and every comparison below hold
with both sides zero in those cases.  The head condition
\texttt{IsLocPt} alone does not impose positive Gram.  In proving
the nonzero branch we now assume positive Gram and choose an
actual representative, as in the definition of the orbital vector.

We first put the actual point into a normal form that preserves its
right Jacquet value.  Write $V=E\oplus W\oplus E^*$ in paired
coordinates, so that
\[
 J=\begin{pmatrix}0&0&I\\0&J_W&0\\I&0&0\end{pmatrix},
 \qquad E=\langle e_0,\ldots,e_{r-1}\rangle.
\]
For a point $g$ over $a$, put $y_b=g^{-1}e_b=(a_b,w_b,b_b)$.
The paired matrix is
$Z_a=\operatorname{diag}(H_h,tJ_s)$, with $J_s$ the reversing
matrix.  The source convention gives
\[
       (b_b)_c=\overline{(Z_a)_{bc}},\qquad
       h(z',w_b)=(Z_a)_{b,r}.
\]
Consequently $b_p=0$, and the $b_b$ with $b\ne p$ form a basis
of the coordinate hyperplane $B=\{b_p=0\}\subset E^*$.
Indeed their square submatrix is the conjugate of
$\operatorname{diag}(H_h,tJ_{s-2})$.  Moreover
$h(z',w_b)=0$ for $b\ne p$, whereas $h(z',w_p)=t$.
The vectors $y_b$ are pairwise isotropic.  In particular $w_p$
is isotropic, while $h(z',z')\ne0$ and $t\ne0$.

Choose $X:E^*\to W$ on $B$ by $Xb_b=w_b$ for $b\ne p$.
Its image there lies in $z'^\perp$.  Isotropy gives
\[
        a_p+X^*J_Ww_p\in\mathbb C e_p.
\]
The projection of $w_p$ to $z'^\perp$ has nonzero norm:
its norm is $-|h(z',w_p)|^2/h(z',z')$.  Thus the remaining
column $Xe_p\in z'^\perp$ can be chosen so that
$a_p=-X^*J_Ww_p$.  Prescribe a skew-Hermitian matrix $K$ on $B$ by
\[
             Kb_b=a_b+\tfrac12X^*J_WXb_b \quad(b\ne p).
\]
This prescription extends to all of $E^*$: on $B\times B$
the skew-Hermitian compatibility is precisely
$h(y_b,y_c)=0$; the missing column is then the negative conjugate
of the already prescribed corresponding row, and the remaining
diagonal entry can be any purely imaginary number.  Direct block
multiplication shows that
\begin{equation}\label{terminal-degree-n0}
 n_0=\begin{pmatrix}
 I&-X^*J_W&K-\tfrac12X^*J_WX\\0&I&X\\0&0&I
 \end{pmatrix}
\end{equation}
is unitary, fixes $E$ pointwise and fixes $z'$, and
\[
 n_0^{-1}y_b=(0,0,b_b)\ (b\ne p),\qquad
 n_0^{-1}y_p=(0,w_p,0).
\]
It is in the original group $N_0(z')$.  Its flag-root coordinates
are zero, and its last character coordinate is zero because
$X^*J_Wz'=0$.  Hence its original character is exactly one.
This proves the required normal form without dividing by any
internal minor of the last block \cite{terminal-degree-source}.

Let $\mathcal RH$ be the original right Jacquet average used in
the inverse-parameter theorem.  It obeys
\begin{equation}\label{terminal-degree-middle-character}
 E_{e_b,e_c}\mathcal RH=\mathbf1_{c=b+1}\mathcal RH,
 \qquad
 E_{e_b,q}\mathcal RH=\mathbf1_{b=r-1}z'_q\mathcal RH
                   \quad(q\in W).
\end{equation}
These are the derivatives of the prescribed character with
$c_{\rm Std}=-2i$, including its middle coordinates.  They apply
to a rightmost generator in any ordered differential word.
The contracted minor commutes with this right averaging by its
actual covariance; the compact support makes differentiation under
the integral valid.  Its covariance under $n_0$ and
$\psi_{z'}(n_0)=1$ permit the calculation at $gn_0$.

At that point the coefficient rows $\xi_b(q)=h(b_q,n_0^{-1}y_b)$
are supported on $E\setminus e_p$ if $b\ne p$, and on $W$ if
$b=p$.  Thus the only nonzero column sets in the contracted
size-$r$ minor are
$(e_0,\ldots,\widehat e_p,\ldots,e_{r-1},q)$ with $q\in W$.
Writing $\epsilon_m=(-1)^{m(m-1)/2}$, their coefficients are
\[
  (-1)^{p+r-1}\det H_h\,t^{s-2}\epsilon_{s-2}
                                      h(b_q,w_p).
\]
Expand the quantum minor in its last column.  The rightmost
middle generator in \eqref{terminal-degree-middle-character}
kills every row except $e_{r-1}$, whose cofactor sign is one.
Summing over $q$ yields exactly
\[
               \sum_{q\in W}z'_q h(b_q,w_p)=h(z',w_p)=t.
\]
There is no conjugation of $t$ in this formula.  The remaining
upper-chain minor has rows $e_0,\ldots,e_{r-2}$ and columns
$E\setminus e_p$.  Successive last-column expansions, as in
\eqref{off-diagonal-col}, remove its last $s-2$ rows and columns.
Every surviving character and cofactor sign is one; the spectral
parameter stays $u$.  The result is the principal $p$ by $p$
minor $Q_p(u)$.  Since
\[
 (-1)^{p+r-1}\epsilon_{s-2}=\epsilon_{s-1},\qquad
 D_R=\det H_h\,t^s\epsilon_s,
\]
we obtain, at every such point,
\begin{equation}\label{terminal-degree-reduction}
       D_R^{-1}\operatorname{AA}_r(u)\mathcal RH
                       =\frac{(-1)^{s-1}}{t}\,Q_p(u)\mathcal RH.
\end{equation}
The principal minor is monic of degree $p$ in $u$, including
$Q_0=1$.  The shift $u=y+(2r+k)/2$ leaves its leading coefficient
unchanged.

The exact head quotient formula expresses $\mathcal V_H(a)$
as the remaining left quotient integral of $\mathcal RH$.
Every point in that integral has the same $Z_a$ and $D_R$;
the preceding argument therefore applies throughout it.
Compact-family properness gives absolute convergence with the
original measures and original normalized quotient cut.
Polynomial coefficients are integrable separately: evaluate at
$p+1$ distinct parameters and use finite Vandermonde inversion.
Integrating \eqref{terminal-degree-reduction} proves
\eqref{terminal-degree-leading}.

We next verify the radial identity needed to use this degree
decrease with the prescribed word operators.  On the regular
cell write $g=w(c)\gamma$, $c_j\ne0$, and keep the original gauge
\[
 \delta(c)=\prod_{j=0}^{r-1}|c_j|^{2(r-j)+k},\qquad
 p_j=c_j\partial_{c_j},\qquad
 \rho_j=r-j+k/2.
\]
The paired coefficient rows on this cell have only their
diagonal $E$ entries $c_j$.  Thus
$\operatorname{AA}_r(u)=(\prod c_j)Q_r(u)$ there.
For any left-character-equivariant smooth coefficient, and for
its right derivatives, the exact cell root and torus rules are
\begin{align*}
 E_{e_a,e_b}K(w(c)\gamma)
   &=\mathbf1_{a=b+1}\frac{c_a}{c_b}K(w(c)\gamma)
                                             &&(a>b),\\
 E_{e_j,e_j}K(w(c)\gamma)
   &=c_j\partial_{c_j}\bigl(K(w(c)\gamma)\bigr).
\end{align*}
For the first rule, conjugating the lower root by the cell
matrix turns it into the left upper root, multiplied by
$c_a/c_b$; the original character is zero unless $a=b+1$.
For the second, the two real diagonal curves multiply $c_j$
by $e^t$ and $e^{it}$, respectively; their complex combination
is $c_j\partial_{c_j}$.  These calculations retain the actual
cell matrix, including its reversed paired order.

For completeness, the ordered determinant recurrence uses a
last-column expansion, followed on its sole surviving nonprincipal
term by a first-factor row expansion.  All far upper roots vanish
by the right rule, and all far lower roots by the left rule.
The adjacent term has the negative cofactor sign and coefficient
$c_m/c_{m-1}$.  More explicitly, that nonprincipal minor has rows
$(0,\ldots,m-2,m)$ and columns $(0,\ldots,m-1)$.
Rotate its last row to first, with sign $(-1)^{m-1}$, and expand
that first row with its generator on the left.  Only column
$m-1$ survives the left rule, with the same sign $(-1)^{m-1}$,
and the remaining minor is $Q_{m-1}(u)$.  Commuting generators
in this row expansion uses
$[E_{ab},E_{cd}]=\delta_{bc}E_{ad}-\delta_{da}E_{cb}$;
the commutator terms cancel the change in the positional shifts.
Equivalently, for the evaluation functional $\mathcal L$ and
any product $P$ of diagonal generators with indices at least
$m+1$, the resulting identity is
\begin{align}\label{terminal-degree-toda}
 \mathcal L(Q_{m+1}(u)P)
   ={}&\mathcal L\bigl(Q_m(u)(u-m+E_{e_m,e_m})P\bigr)\notag\\
    &-\frac{c_m}{c_{m-1}}\mathcal L(Q_{m-1}(u)P)
                                                     \quad(m\ge1).
\end{align}
Here is the all-rank first-factor formula used above, for arbitrary
ordered injective row and column lists of length $q$:
\[
 Q_{a,b}(u)=\sum_{j=0}^{q-1}(-1)^j
 T_{a_0b_j}(u-q+1)
 Q_{a\setminus a_0,\,b\setminus b_j}(u).
\]
To prove it, write the defining column alternant and group terms
in pairs that exchange two adjacent assigned rows.  The displayed
two-factor identity in the inverse-parameter proof replaces each
such paired alternant by the corresponding row-ordered alternant,
reversing the two positional shifts.  Repeating at the next
adjacent position moves the chosen row past one more factor;
induction on the number of positions gives the row alternant with
parameters $u-q+1,u-q+2,\ldots,u$.  Group that alternant by the
column assigned to its first row.  Its sign is $(-1)^j$ and the
remaining alternant has parameters $u-q+2,\ldots,u$, which is
the length-$(q-1)$ minor at parameter $u$.  This proves the formula
for every $q$, starting with $q=1$, and justifies its use above.
It explains both
the unshifted $u$ in the two smaller minors and the $u-m$ in
the diagonal factor.  The auxiliary diagonal factors commute
with the root generators being removed, so induction retains
\eqref{terminal-degree-toda} with those factors present.

The gauge identity
\[
       p_j\delta=\rho_j\delta,\qquad
       \delta\,(y+\rho_j+p_j)K=(y+p_j)(\delta K)
\]
now gives, by induction from $Q_0=1$ and $Q_1(u)=u+E_{00}$,
\begin{equation}\label{terminal-degree-vab}
 A_0=1,\quad B_0=0,\qquad
 A_{j+1}=(y+p_j)A_j-c_jB_j,\quad B_{j+1}=c_j^{-1}A_j.
\end{equation}
These operators act on the actual gauged Jacquet image.
Right derivatives preserve left covariance,
and compact properness justifies all differentiations.  Thus this
is the original radial comparison, not a replacement differential
model.  The original terminal coordinate satisfies
$D_R=(\prod c_j)\zeta$, so division by $D_R$ turns
\eqref{terminal-degree-vab} into precisely
$\zeta^{-1}A_r=\texttt{TL.aP}(y)$.  All original Haar constants
occur on both sides of the same Jacquet integral and remain
unchanged.  Also $\mathcal A_y$ preserves compact support and
the positive Gram cutoff.  This proves \texttt{RadAATop}
\cite{terminal-degree-radial-source}.

Here is the exact first-component word comparison.  Write
$\mathcal F_w$ for the prescribed terminal frame and
$\mathcal B_w$ for its ordered site product.  The plus update is
\[
 \mathcal F_{(+y,w)}=
 \iota_{\rm new}\!\left[
 \begin{pmatrix}\zeta^{-1}&0\\0&\bar\zeta^{-1}\end{pmatrix}
 \begin{pmatrix}1&0\\G\ell_w(y)&G\end{pmatrix}\right]
 \mathcal F_w.
\]
The old frame is free of the new tensor position; its entries
commute with the new site operators.  Move it past the new site
product, whose first column is $(A_r,B_r)^t$, and multiply the
displayed two by two matrices.  Their first component is
$\zeta^{-1}A_r$.  Therefore the new bit-zero component is
\texttt{TL.aP}(y) applied to the old component.  All old frame
entries, rational inverses and word factors remain in the old
component; none has been omitted or bounded.  This is the
bit-zero part of the actual \texttt{WordStep} identity
\cite{terminal-degree-word-source}.

The differential realization supplies a smooth compact
representative $H_{w,e}$ of $\Theta_{w,e}F$ on the positive
Gram domain.  Indeed each of its finitely many terms is a
finite-order local differential operator there, so it preserves
the compact support and vanishes below the same positive Gram
cutoff; extension by zero is smooth.  The original orbital
vector agrees with that representative at every head, by the
same quotient formula.  The regular radial comparison just
proved and the first-component update show that
$H_{(+y,w),(0,e)}$ and $\mathcal A_yH_{w,e}$ have identical
regular gauged Jacquet images.  The already proved every-head
density theorem applies to their compact difference and gives
\[
     \mathcal V_{\Theta_{(+y,w),(0,e)}F}(a')
                       =\mathcal V_{\mathcal A_yH_{w,e}}(a')
                                   \quad\hbox{for every head }a'.
\]
This also proves the full \texttt{TopFirst}, for any compact
representative agreeing on the stated positive domain.

Finally apply leading-coefficient interpolation to
\eqref{terminal-degree-leading} repeatedly.  Process the new
letters in reverse list order, because the rightmost operator
acts first.  Each extraction multiplies the old orbital vector
by $\lambda$.  Every intermediate word is a suffix of one of
the assumed admissible full words and is therefore admissible.
Reversing the grid indices back gives exactly
\eqref{terminal-degree-grid}; scalar Lagrange weights commute,
so this reversal changes neither their product nor the sum.
For $p=0$ each grid is a singleton and its weight is one.  For
$K=0$ both sides are the same old orbital vector.  If $r=0$
there is no terminal head with $s\ge2$, so the degree assertion
has no instance.  All these cases retain the original conventions.
\end{proof}
